\documentclass[12pt]{article}
\usepackage[T1]{fontenc}
\usepackage[utf8]{inputenc}
\usepackage{lmodern}
\usepackage{textcomp}
\usepackage{enumitem}
\usepackage{geometry}
\usepackage{graphicx}
\usepackage{amsmath, amssymb}
\geometry{margin=1in}
\begin{document}
\begin{center}
Philosophy - Graduate Level Assessment 25 \
Total Marks: 40 \
Answer all questions.
\end{center}

\section*{Multiple Choice (10 marks)}
\begin{enumerate}[label=\arabic*.]
\item " \_Ad lazarum\_ " is a specific kind of
\begin{enumerate}[label=(\alph*)]
    \item Fallacy of composition
    \item Complex cause
    \item Red herring
    \item False sign
    \item Hasty generalization
\end{enumerate}
\item Logical behaviorism is at odds with our commonsense intuition that mental states \_\_\_\_\_.
\begin{enumerate}[label=(\alph*)]
    \item cause behavior
    \item are the same as behavioral states
    \item exist
    \item can be directly observed
    \item are always visible
\end{enumerate}
\item Van den Haag is
\begin{enumerate}[label=(\alph*)]
    \item a capitalist.
    \item a feminist.
    \item an anarchist.
    \item a socialist.
    \item a distributionist.
\end{enumerate}
\item In his discussion of discrimination in war, Valls suggests that
\begin{enumerate}[label=(\alph*)]
    \item the concept of combatants and noncombatants is outdated and irrelevant.
    \item "terrorism" should be defined as violence against noncombatants.
    \item there is no difference between combatants and noncombatants.
    \item the difference between combatants and noncombatants is largely ambiguous.
    \item the difference between combatants and noncombatants is categorical and clear.
\end{enumerate}
\item According to Sinnott-Armstrong, the fact your government morally ought to do something
\begin{enumerate}[label=(\alph*)]
    \item does not prove that government officials ought to promote it.
    \item does not prove that you ought to do it.
    \item proves that you ought not to do it.
    \item proves that you ought to do it, too.
\end{enumerate}
\end{enumerate}

\section*{True / False (10 marks)}
\textit{Answer True or False}
\begin{enumerate}[label=\arabic*.]
\setcounter{enumi}{5}
\item True or False: Moore argues that many past philosophers have committed the naturalistic fallacy in their attempts to define 'good.'
\item True or False: Mill argues that the distinction between justice and other moral obligations corresponds to the difference between positive and negative duties.
\item True or False: Baxter posits that introducing new species is essential for achieving ecological balance in addressing environmental issues.
\item True or False: Metz contends that dignity is primarily based on a capacity for scientific understanding.
\item True or False: Macedo does not identify cosmopolitan egalitarian duties as a societal obligation toward nonmembers.
\end{enumerate}

\section*{Long Form Response (20 marks)}
\textit{Answer each question with a clear and complete explanation.}
\begin{enumerate}[label=\arabic*.]
\setcounter{enumi}{10}
\item Discuss the role of the minor premise in determining the validity of a red herring argument, analyzing how it affirms the antecedent or denies the consequent.
\item Using indirect truth tables, evaluate the validity of the argument H ≡ (\textasciitilde{}I ∨ J), H ∨ \textasciitilde{}J / \textasciitilde{}I. If invalid, provide a counterexample and analyze its implications.
\end{enumerate}

\vfill
\noindent\textit{End of Paper}
\end{document}